\section*{Conventions and notations}\label{NOTATIONS}
\addcontentsline{toc}{section}{Conventions and notations}
Throughout this thesis, we let $n$ be a positive integer. We usually let $\cat{A}$ denote an additive category. When not stated otherwise, the word kernel (respectively cokernel) will refer to the kernel morphism (respectively cokernel morphism) instead of the object. For a morphism $f : A \xra{} B$, we call $A$ its domain and $B$ its codomain. For two objects $A, B \in \cat{A}$, we let $\cat{A}(A,B)$ denote the set of morphisms from $A$ to $B$. We will use $\ker{f}$ to refer to its kernel morphism and $\cok{f}$ for its cokernel morphism, whereas $\Ker{f}$ will refer to its kernel object and $\Cok{f}$ its cokernel object. For any object $A \in \cat{A}$, we denote by $1_A$ the identity morphism on $A$. We let $\Chlen{n}{\cat{A}}$ denote the set of chain complexes on $\cat{A}$ concentrated in degrees $\{0, \ldots, n\}$ and $\Ch{\cat{A}}$ denotes the set of all chain complexes on $\cat{A}$. We let $\text{H}(\cat{A})$ denote the homotopy category of $\cat{A}$


